% Copyright 2024- Jay Jay Billings. Some rights reserved.
\documentclass{article}
\usepackage{graphicx} % Required for inserting images
\usepackage{hyperref}
\usepackage{tabularx}

\title{Stabbing the ASUS ME571K\_SUB rev1.4}
\author{Jay Jay Billings, Ph.D.}
\date{\today}

\begin{document}

\maketitle

\section*{Voltages}

**INSERT** Picture of the board with breakout board and pin layout.

\begin{table}[h!]
\centering
\begin{tabularx}{ 0.8\textwidth }{ | >{\raggedright\arraybackslash}X | >{\centering\arraybackslash}X | }
 \hline
 \multicolumn{2}{|c|}{Pin Voltages} \\
 \hline
 Voltage & Pin Number(s) \\
 G & 6, 9, 18, 23, 28, 35, 36, 38, 40 \\
 0V with $R\approx350\Omega$, 31, 39 \\
 5V & 10, 11, 12, 13, 14, 15, 16, 17, 18 \\
 0.2V-1.6V oscillating & 29, 30, 42 \\
 \hline
\end{tabularx}
\caption{Notable pin voltages for the board. G denotes ground. All unlisted pins had nominal stable voltages $\approx100$mv.}
\end{table}

\end{document}
